blé d'Espagne » au chapitre de Sainte Cécile d'Albi représenté par
Guillaume Lavergne, vicaire perpétuel (curé) de Négremont. En 1678,
les habitants d'Alban menaceront même d'émigrer. Et en 1682, il n'y a
que 21 habitants pour payer l'impôt dit de la « pesade ». À cette date,
le cheptel d'Alban ne comprend que 6 vaches, 4 bœufs, 125 bêtes « à
laine »... c'est peu !

Le nom d'un curé de cette époque nous est révélé par un acte signé de
lui. C'est Louis de Ménard qui est en 1658 le curé d'Alban. Le nom de
deux consuls pour l'année 1669 est venu jusqu'à nous : il s'agit de
Jacques Calmet et Jean Capel. Dans le même temps, un Pierre Roussel
est curé de Rayssac et son neveu en assurera la charge en 1685.
(Ces Roussel sont originaires de Saint André).

On peut noter en souvenir de celle qui fut « Madame de la Tonié »,
si connue à Alban jusqu'en 1948, date de sa mort, qu'un décret de
Colbert, du 24 mai 1667, ayant ordonné la révision des titres de
noblesse, René de Clergue\footnote{?} de la Tonié, put prouver ses titres en faisant
valoir ses six quartiers de noblesse et justifier ainsi la valeur de son
blason, qui héraldiquement se lit ainsi : « de gueule à deux pans et
deux bandes d'or portant de chef trois essans\footnote{Blason Alban?}, l'or, le tout surmonté
d'un casque d'argent empanaché d'or et de gueule (= sang et or) ».

Angèle de la Tonié était née Terral et ce pouvait être, de ce fait,
la tante et la marraine de Denise Aigrot, propriétaire du Chalet au
Noyer.